\section*{Exercises about Diagonalizable Operators}

\exercise{2}
\begin{xrcs}
  Suppose $T \in \linmap(V)$ has a diagonal matrix $A=\m{T}$ with respect to some basis of $V$. Prove that if $\lambda \in \myF$, then $\lambda$ appears on the diagonal of $A$ precisely $\dim E (\lambda , T)$ times.

  \begin{xprf}
    If $\lambda$ is not an eigenvalue of $T$, then
    \begin{equation}
      \dim E(\lambda, T) = \dim \{v \in V \mid Tv = \lambda v\} = \dim \{0\} = 0,
    \end{equation}
    and therefore $\lambda$ does not appear among the diagonal entries of $A$.

    Now assume $\lambda$ is an eigenvalue. Let $\mathcal{B} = \{v_1, \dots, v_\ell\}$ be a basis of $V$ with respect to which $A = \m{T}$ is diagonal; write
    \begin{equation}
      A = \operatorname{diag}(\mu_1, \dots, \mu_\ell),
    \end{equation}
    so that for each $j$ we have $Tv_j = \mu_j v_j$. In particular, every $v_j$ is an eigenvector of $T$.

    Define the subset of basis vectors corresponding to the eigenvalue $\lambda$ by
    \begin{equation}
      S = \{\, v_j \in \mathcal{B} \mid \mu_j = \lambda \,\}.
    \end{equation}
    We first observe that $\myspan S \subseteq E(\lambda, T)$, since each vector in $S$ is by definition a $\lambda$-eigenvector and linear combinations of $\lambda$-eigenvectors remain in the $\lambda$-eigenspace.

    For the reverse inclusion, let $w \in E(\lambda, T)$. Express $w$ in the basis $\mathcal{B}$:
    \begin{equation}
      w = \sum_{j=1}^\ell c_j v_j \qquad (c_j \in \myF).
    \end{equation}
    Applying $T$ to both sides and using $Tv_j = \mu_j v_j$ gives
    \begin{equation}
      \lambda \sum_{j=1}^\ell c_j v_j = Tw = \sum_{j=1}^\ell c_j \mu_j v_j.
    \end{equation}
    By uniqueness of coordinates with respect to the basis $\mathcal{B}$, the coefficient of $v_j$ on the left must equal the coefficient of $v_j$ on the right for each $j$, hence
    \begin{equation}
      (\lambda - \mu_j) c_j = 0 \quad \text{for every } j = 1, \dots, \ell.
    \end{equation}
    Consequently, $c_j = 0$ whenever $\mu_j \neq \lambda$, so every nonzero coefficient $c_j$ corresponds to some $v_j \in S$. Thus $w$ is a linear combination of vectors from $S$, meaning $w \in \myspan S$. This proves the reverse inclusion and therefore
    \begin{equation}
      \myspan S = E(\lambda, T).
    \end{equation}

    Since $S$ is a subset of the basis $\mathcal{B}$, its vectors are linearly independent; hence the number of elements of $S$ equals the dimension of its span:
    \begin{equation}
      |S| = \dim \myspan S = \dim E(\lambda, T).
    \end{equation}
    Finally, by definition $|S|$ is exactly the number of diagonal entries of $A$ equal to $\lambda$. We conclude that the scalar $\lambda$ appears on the diagonal of $A$ precisely $\dim E(\lambda, T)$ times.
  \end{xprf}

  \begin{failedattempt}
    If $\lambda$ is not an eigenvalue, then $\dim E(\lambda, T) = \dim \set{ v : Tv = \lambda v } = \dim \set{0} = 0$, and it does not appear on the diagonal of $A$.
    Now, assume $\lambda$ is an eigenvalue of $T$. Since $T$ has a diagonal matrix with respect to some basis of $V$, say $v_1, \ddd, v_\ell$ with $\ell = \dim V$,
    all those vectors are eigenvectors of $T$. A subcollection that satisfies $Tv = \lambda v$ forms a basis of $E(\lambda, T)$.
    Let $v_{n_1}, \ddd, v_{n_m}$ be a basis of $E(\lambda, T)$ with $\dim E(\lambda, T) = m$ and $n_j \in \set{1, \ddd, \ell}$ for $j = 1, \ddd, m$.
    This implies that $T v_{n_j} = \lambda v_{n_j}$ for $j \in \set{1, \ddd, m }$, which can only happen if $\lambda$ appears on $A$ exactly $m = \dim E(\lambda, T)$ times.
  \end{failedattempt}

  \begin{ainote}[Where it breaks]
    One sentence carries the entire exercise and is asserted rather than proved: \qt{A subcollection that satisfies $Tv = \lambda v$ forms a basis of $E(\lambda, T)$.}

    Half of it is easy. Those basis vectors do lie in $E(\lambda, T)$, and being part of a basis they are linearly independent. What is \emph{not} obvious is that they \emph{span} $E(\lambda, T)$ --- i.e. that the eigenspace contains nothing beyond what those particular basis vectors generate. Nothing said so far rules out some $w \in E(\lambda,T)$ having a nonzero coefficient on a $v_j$ whose eigenvalue $\mu_j$ differs from $\lambda$, and without that ruled out the count $|S| = \dim E(\lambda,T)$ is unsupported.

    Closing the gap is the whole of the reverse inclusion in the corrected proof: expand $w$ in the basis, apply $T$, and compare coordinates to get $(\lambda - \mu_j)c_j = 0$, which forces $c_j = 0$ exactly when $\mu_j \neq \lambda$. The step that makes it work is \emph{uniqueness} of coordinates with respect to a basis --- the one property that distinguishes a basis from a mere spanning list.

    The lesson generalises: whenever an argument says \qt{take the subcollection with property $P$ --- it is a basis of the relevant subspace}, independence is usually free and spanning is where the work is.
  \end{ainote}
\end{xrcs}

\exercise{14}
\begin{xrcs}
  Suppose $S, T \in \linmap(V)$ are diagonalizable operators on a finite-dimensional vector space $V$ that commute ($ST = TS$). Prove that there exists a basis of $V$ consisting of vectors that are eigenvectors of both $S$ and $T$.
\begin{xprf}
  Since $S$ is diagonalizable, $V$ decomposes as a direct sum of the distinct eigenspaces of $S$:
  \[
    V = E(\lambda_1, S) \oplus \dots \oplus E(\lambda_m, S).
  \]
  Because $S$ and $T$ commute, each eigenspace $E(\lambda_j, S)$ is invariant under $T$ (by Exercise 5E.2).
  Since $T$ is diagonalizable, its restriction $T|_{E(\lambda_j, S)}$ to the invariant subspace $E(\lambda_j, S)$ is also diagonalizable.
  Therefore, for each $j \in \{1, \dots, m\}$, there exists a basis $\mathcal{B}_j$ of $E(\lambda_j, S)$ consisting of eigenvectors of $T$.
  Every vector in $\mathcal{B}_j$ is an eigenvector of $T$, and (since it belongs to $E(\lambda_j, S)$) it is also an eigenvector of $S$ with eigenvalue $\lambda_j$.
  The union $\mathcal{B} = \mathcal{B}_1 \cup \dots \cup \mathcal{B}_m$ is a basis of $V = \bigoplus_{j=1}^m E(\lambda_j, S)$ consisting of vectors that are simultaneous eigenvectors of both $S$ and $T$.
\end{xprf}
\end{xrcs}